\newpage
\section{Misión 5: El dinero es la raíz}\label{sec:mission5}

\begin{minipage}{0.90\linewidth}
Toda la colonia bulle de energía intelectual tras instalar la HugeAntenna\texttrademark\
y conectarla a Internet de la Tierra -- ¡piensa en todo lo bueno que se han perdido hasta ahora!
%
Por desgracia, ahora que la colonia está enganchada a Internet, el tiránico gobierno de la Tierra ha decidido encarecer mucho el uso del enlace de la antena interplanetaria.
Los precios abusivos causan graves problemas en la sociedad de la colonia, demostrando otra vez que el dinero \mbox{-- o,} quizá, su \mbox{falta --} puede ser la raíz de todos los males.
%ñ
¿Ayudarás a la colonia a recuperar la estabilidad financiera?

\end{minipage}
\hspace{0.015\linewidth}
\begin{minipage}{0.08\linewidth}
% \vspace{-0.5cm}

\includegraphics[width=\linewidth]{./cash.jpg}
\end{minipage}

% \vspace{0.3cm}
\noindent
\textbf{\raisebox{0.03cm}{\warning}¡Atención! Esta misión es crítica}. Tómate tu tiempo, pide toda la ayuda que necesites y mantén el ánimo. Las misiones anteriores pueden considerarse un entrenamiento para esta -- se recomienda encarecidamente terminarlas con cuidado antes de abordarla.
\vspace{-0.3cm}

\subsection*{Detalles del informe}

La Tierra acepta negociar el precio en Mar\$ (\mbox{Dollaroos marcianos}) por byte transmitido por la HugeAntenna\texttrademark.
El precio depende de la MTU de la antena, según la figura inferior.
Observa que $\textrm{MTU} \in \left\{200,\,400,\,\ldots,\,3000\right\}$.
La asamblea necesita un análisis cuantitativo que justifique la mejor elección de MTU.

\centerline{%
\begin{minipage}{0.5\linewidth}
\includegraphics[width=\linewidth]{cost_plot.pdf}
\end{minipage}
}

\noindent
Para simplificar, considera un escenario representativo en el que un usuario de una base de la colonia envía
\textit{payloads} (carga útil) de exactamente $1900$~bytes directamente sobre~IP.

\begin{enumerate}[itemsep=0.1cm]
\item Da un mapa topológico de la red de la colonia hasta el enlace entre la HugeAntenna\texttrademark\  y la Tierra, sin incluirlo. En esas LANs (red de área local), indica la MTU más pequeña que evita la fragmentación dentro de la colonia.
Puedes reutilizar la topología de misiones anteriores o crear una nueva.

\item Considera al menos $4$ valores posibles para la MTU de la HugeAntenna, $2$ menores que $1900$~bytes y dos mayores. Para cada valor de $\textrm{MTU}_\textrm{HugeAntenna}$, indica:
  \begin{itemize}
  \item Cuántos paquetes se envían por la Huge Antenna para transmitir el payload.
  \item Qué campos no son idénticos en esos paquetes (solo si se envía más de un paquete).
  \item El coste total de enviar el payload.
  \end{itemize}
Coordínate con los demás ingenieros para hallar el valor óptimo global de $\textrm{MTU}_\textrm{HugeAntenna}$.

\item Compara el tráfico total (en bytes) transmitido dentro de la colonia al enviar el primer payload (con todas las cachés vacías) frente al segundo (con las cachés llenas y válidas).
\end{enumerate}


\subsection*{Detalles de la retroalimentación}

Sigue un proceso de retroalimentación similar al de misiones anteriores, asegurándote de abordar:
\begin{itemize}
\item ¿Están las MTU de la colonia configuradas según lo pedido?
\item ¿Se considera correctamente la fragmentación al determinar el coste?
\item ¿Se contabiliza bien la sobrecarga debida a ARP al calcular el tráfico total?
\end{itemize}

